\documentclass[11pt,a4paper]{article}
\usepackage[T1]{fontenc}
\usepackage[utf8]{inputenc}
\usepackage{amsmath,amssymb,amsthm}
\usepackage[margin=1in]{geometry}
\usepackage[colorlinks=true,linkcolor=blue,urlcolor=blue]{hyperref}
\theoremstyle{plain}
\newtheorem{theorem}{Theorem}
\newtheorem{lemma}[theorem]{Lemma}
\newtheorem{proposition}[theorem]{Proposition}
\newtheorem{corollary}[theorem]{Corollary}
\theoremstyle{definition}
\newtheorem{remark}[theorem]{Remark}
% A quoted conjecture can be a single hundred-term formula whose terms are thirty-digit
% numbers. TeX cannot hyphenate those, so without a little stretch every such quote runs off
% the page; the linked-a-file papers are all of that shape.
\setlength{\emergencystretch}{3em}

\title{The empirical recurrence for OEIS A314015, proved from an exact model}
\author{Adrian Perez Fontelles\\ \small Independent researcher}
\date{14 September 2026}
\begin{document}
\maketitle

\begin{abstract}
OEIS A314015 carries an empirical recurrence of order $6$ contributed by Chai Wah Wu. It is true, and it is
decidable rather than empirical. The entry's sequence is given exactly by a certified distance function on the tiling, counted by Ehrhart, from
which $a$ provably satisfies a MONIC linear recurrence of order $S=12$, derived below and not
assumed. The residual of the conjectured recurrence satisfies the same one, so whether it
annihilates $a$ is settled by a finite exact computation. No transfer matrix is involved and the count is not a walk.
\end{abstract}

\noindent\small 2020 Mathematics Subject Classification. 52C20, 05B45, 52C07, 05A15.\normalsize

\section{The conjecture}

OEIS A314015 is ``Coordination sequence Gal.4.100.3 where Gal.u.t.v denotes the coordination sequence for a vertex of type v in tiling number t in the Galebach list of u-uniform tilings.''. It has offset $0$ and begins
\[
1, 5, 11, 15, 19, 25, 30, 35,\ \dots
\]
The entry states, as a conjecture contributed by its author and never marked settled:
\begin{quote}
Conjectures from \_Chai Wah Wu\_, Apr 17 2024: (Start) a(n) = 2*a(n-1) - 2*a(n-2) + 2*a(n-3) - 2*a(n-4) + 2*a(n-5) - a(n-6) for n > 6.
\end{quote}
As of the ``Last modified'' line on the live entry (Apr 18 2024 17:56:45, revision 7)
this is still recorded as empirical, and nothing on the entry records it as proved.

\section{The tiling, the lattice, and the distance}

There is no array here and no transfer matrix. The entry names a vertex of one of Brian
Galebach's $k$-uniform tilings, and that tiling is rebuilt exactly: the expanded notation
fixes, for every vertex type, the angles between its edges and the type and frame of the
vertex at the far end of each. All angles are multiples of $15^\circ$, so every vertex
position is a $\mathbb{Z}$-combination of $24$th roots of unity and lives in
$\mathbb{Z}[\zeta_{24}]=\mathbb{Z}^{8}$ modulo $x^{8}-x^{4}+1$. Vertex identity is
therefore decided by integer equality and not by a tolerance.

Let $L$ be the group of translations of the plane carrying the tiling to itself. Two
independent generators are found and each is VERIFIED to be a symmetry: for every vertex $w$
of an inner patch, $w+t$ is a vertex of the same type and the neighbours of $w+t$ are exactly
the neighbours of $w$ translated. (Preserving a vertex's type and its set of edge directions
is NOT sufficient, because a vertex figure with a rotational symmetry has several frames with
the same direction set.) The vertices fall into $30$ orbits under $L$, and a vertex is
written $(c,m)$ with $c$ its orbit and $m\in\mathbb{Z}^2$ its coordinates in a basis of $L$.

Write $D_c(m)$ for the graph distance from the entry's own vertex to $(c,m)$. The claim proved
below is that on each orbit $D_c$ is a maximum of finitely many affine functions,
\[
D_c(m)\;=\;\max_i\bigl(A_{c,i}m_1+B_{c,i}m_2+C_{c,i}\bigr).
\]

\section{Why that is the distance, and not merely a fit}

The affine pieces are obtained from a breadth-first patch, but nothing below rests on that. A
function $D$ with $D(\text{base})=0$ satisfying

\begin{itemize}
\item[(b)] $D_{c'}(m+d)\le D_c(m)+1$ for every edge $(c,c',d)$ of the quotient and every $m$;
\item[(c)] every vertex other than the base has an edge attaining $D_c(m)=D_{c'}(m+d)+1$;
\end{itemize}

IS the distance. (b) gives $D\le d$ by induction along a shortest path, and (c) gives $D\ge d$
by induction on $D$; neither induction cares where $D$ came from. So the fit is a candidate and
these two conditions are the proof.

Both are decided exactly, and on REGIONS rather than on sampled points. The set where the
$i$-th piece is the maximum is the polyhedron
$P_{c,i}=\{m: \ell_{c,i}(m)-\ell_{c,j}(m)\ge0 \text{ for all } j\}$, so (b) is an
affine inequality on a polyhedron --- true exactly when it holds at every vertex of that
polyhedron and its gradient pairs non-negatively with every recession ray. Condition (c) is
not an affine statement but a disjunction, and it says that $P_{c,i}$ is COVERED by the
regions
$S_k=\{m:\ell_{c,i}(m)-1-\ell''_{k,j}(m+d_k)\ge0\text{ for all }j\}$, one per
neighbour $k$; that is decided by subtracting the $S_k$ from $P_{c,i}$ one at a time and
asking whether any region is left. Every coefficient is an integer, and since only lattice
points matter a violated constraint is written $\le-1$, so no strict inequality occurs
anywhere. A bounded box around the base is checked exhaustively and the four half-planes
outside it are covered by the region arithmetic.

\section{The bound}

The ball of radius $t$ meets orbit $c$ in the lattice points of
\[
P_c(t)=\{m\in\mathbb{Z}^2:\;A_{c,i}m_1+B_{c,i}m_2\le t-C_{c,i}\ \text{ for every } i\},
\]
a polygon whose facet NORMALS are fixed and whose right-hand sides move linearly with $t$.
Every candidate vertex of it is the meet of two facets and moves affinely in $t$, so whether a
given vertex satisfies a given facet flips at most once, at one rational $t$; past the largest
such value --- computed here, not assumed --- the combinatorial type never changes again. By
Ehrhart's theorem for such a family the lattice count is then a quasi-polynomial in $t$ of
degree $2$, whose period divides the lcm of the $2\times2$ determinants of pairs of facet
normals, since those are the only denominators a vertex can have. Write $q$ for that lcm;
here $S=2q=12$.

Summing over orbits, $|B(t)|$ is a degree-$2$ quasi-polynomial of period $q$ past the onset,
so $a(n)=|B(n)|-|B(n-1)|$ is quasi-LINEAR of period $q$ there, and so is the residual of any
fixed linear recurrence. A quasi-linear function of period $q$ that vanishes at $2q$
consecutive indices past the onset vanishes at every later one, which is what makes the check
below finite and complete rather than a sample. The counting itself is never done through the
closed form: the lattice points are counted exactly at every radius, and the quasi-polynomial's
only job is to supply $q$.

The annihilator is therefore $A(z)=(z^{q}-1)^2$, of degree $S=2q$ and with $A(0)=1\ne0$, and
it annihilates $a$ from the onset on rather than everywhere: the indices below it are a
genuine transient, exactly as $\dim S_0$ and $\dim S_2$ are for the cusp-form family. That
costs nothing, because the residual below is computed term by term from the exact lattice
count and the run of zeros the theorem uses lies entirely above the onset.

\section{The proof}

\begin{lemma}\label{lem:crit}
Let $A(z)=z^S-\sum_{j<S}\alpha_jz^{j}$ be the monic annihilator derived above and
$q(z)=z^{6}-\sum_ic_iz^{6-i}$ the characteristic polynomial of the conjectured
recurrence. The residual $r(n)=a(n)-\sum_ic_ia(n-i)$ is a linear combination of shifts of $a$,
so $A$ annihilates $r$ as well. Since $A(0)\ne0$ the recurrence it gives runs in both
directions, and $S$ consecutive zeros of $r$ force $r$ to vanish at every later index.
\end{lemma}

\begin{proof}
$A$ annihilates $a$, hence every shift of $a$, hence any linear combination of shifts; that is
$r$. A monic recurrence with nonzero constant term determines each term from the $S$ before it
and each from the $S$ after it, so a run of $S$ zeros propagates.
\end{proof}

\begin{theorem}
\[
a(n)\;=\;2\,a(n-1) - 2\,a(n-2) + 2\,a(n-3) - 2\,a(n-4) + 2\,a(n-5) - a(n-6)
\]
for every $n>6$.
\end{theorem}

\begin{proof}
The terms $a(0),\dots$ were computed exactly in integer arithmetic from the model, extended by
$A$ where the model itself was not evaluated, and $r(n)$ formed for each index. Those integers
vanish from $n=6+1$ onwards, and the run exceeds $S+6$, which by
Lemma~\ref{lem:crit} settles every larger index at once. The bound is exact: at $n=6$ the recurrence fails on the entry's own published terms, so no larger range holds.
\end{proof}

\section{Verification}

Three checks, each independent of the algebra above.

First, the model reproduces all $50$ terms the entry publishes, exactly, in integer
arithmetic. This is what ties the model to the entry: it is built by reading the entry's
English, and a misreading gives different counts.

Second, the conjectured recurrence was evaluated directly on those published terms, with no
model involved, and holds wherever the proved range and the data overlap.

Third, the derived annihilator $A$ was itself tested on terms it had not been given: the model
was evaluated beyond the $S$ values that determine it and $A$ reproduced them. A bound that is
too small fails this test, which is why it is run.

\begin{thebibliography}{9}
\bibitem{oeis} The OEIS Foundation, \emph{The On-Line Encyclopedia of Integer
Sequences}, \url{https://oeis.org}, 2026. Sequence A314015.
\bibitem{beckrobins} M.~Beck and S.~Robins, \emph{Computing the Continuous Discretely},
2nd ed., Springer, 2015. (Ehrhart quasi-polynomials and their periods.)
\bibitem{stanley} R.~P.~Stanley, \emph{Enumerative Combinatorics}, Volume 1, 2nd ed.,
Cambridge University Press, 2011.
\bibitem{ds} F.~Diamond and J.~Shurman, \emph{A First Course in Modular Forms},
Springer GTM 228, 2005.
\end{thebibliography}

\end{document}
